\section{本章小结}
\label{sec:mmaps-summary}

本章开头问：给定流形之间的映射，它怎样对应点，又怎样改变点处的方向和读数？对点而言，光滑映射由两侧坐标图中的坐标表示判断，换用相容的图不改变结论。同一个坐标公式可能描述两件事：同一点在两张图中的记录之间的换算，或者把点实际送到另一个点；前者是恒等映射的坐标表示，后者才移动了点。曲线和实值函数都是光滑映射的特例。

对方向而言，映射的微分$\mathrm dF_p:T_pM\to T_{F(p)}N$把经过$p$的曲线速度推送为经过$F(p)$的曲线速度，在坐标基下的矩阵是坐标表示的Jacobian矩阵；边界点处用光滑延拓的Jacobian矩阵，结果与延拓无关。链式法则$\mathrm d(G\circ F)_p=\mathrm dG_{F(p)}\circ\mathrm dF_p$成立，秩不依赖图，实值函数的微分与第~\ref{chap:manifold-basics}章的余切向量$\mathrm df_p$一致。

对整个向量场逐点推送时，不同原像可能给出同一点处相互冲突的向量，像以外的点没有值，所以一般映射不能把向量场推送成目标上的向量场。能做的是判断两侧的场是否$F$相关，它等价于$X(g\circ F)=(Yg)\circ F$，并与Lie括号相容。若映射是光滑双射且逆映射光滑，推送得到唯一的光滑向量场；逆映射不光滑时，例如$x\mapsto x^3$，推送结果可能不光滑。

对读数而言，余切向量沿相反方向拉回：$(\mathrm dF_p)^*\omega=\omega\circ\mathrm dF_p$，坐标中由Jacobian矩阵的转置给出。协向量场的拉回对任何光滑映射都有定义，因为每个原像点只需要一个目标值；它满足$(G\circ F)^*=F^*\circ G^*$与$F^*(\mathrm df)=\mathrm d(f\circ F)$。推送与拉回的方向不同、作用对象不同，所需的条件也不同。

这些都是映射的局部操作。微分可逆、单射或满射分别说明映射在一点附近保留或覆盖了多少方向，但不能单独决定两个流形整体上是否等价，或一个流形的像能否成为子流形。第~\ref{chap:manifold-relations}章在本章的微分与秩的基础上，加入可逆性与拓扑条件来回答这些问题。
